% Section 1. Introduction.

Behavioral researchers now prompt large language models (LLMs) with demographic personas and treat the answers as \emph{silicon samples}.
\citet{argyle2023outofone} used such samples to reproduce vote choice and opinion by group, \citet{park2024diminished} used them to replicate classic experiments, and \citet{lin2026simulators} describes their use in piloting instruments before fielding them.
A simulated sample costs little, reaches groups that are hard to recruit, and can be rerun whenever the design changes.
Each of these uses assumes that the simulated sample stands in for a real one in the respect the study measures.

That assumption is checked unevenly.
Across eight social science journals, \citet{desai2026validating} found validation practice in LLM studies ``inconsistent and limited.''
\citet{park2024diminished} found GPT-3.5 answering with almost no variation where people vary, and \citet{bisbee2024synthetic} found synthetic opinion distributions narrower than the survey distributions they imitate.
\citet{wang2025misportray} showed that LLMs standing in for participants can misportray and flatten identity groups, and \citet{dominguezolmedo2024questioning} found that models' survey answers trend toward uniform randomness once ordering and labeling biases are adjusted for.
\citet{sarstedt2024silicon} set out guidelines for silicon samples in consumer research, and \citet{lin2026workflow} proposed a workflow in which the evidence a study needs rises with its claim, setting correspondence with a population's means, variance, scale structure and correlations as targets without a test attached to each.

One persona from the corpus analyzed here shows what reading individual answers misses.
Prompted as a single Black woman on Medicaid with a household income under \$35,000, GPT-4o-mini returns a PHQ-8 total of 8 on its first draw, a plausible answer in the mild band. % R: 85_worked_case.csv
Over thirty draws the same prompt averages 10.3, and adults in the National Health and Nutrition Examination Survey (NHANES) who share her race, sex, income and marital status average 4.7. % R: 85_worked_case.csv
None of her single answers shows that the simulated population she belongs to sits well above the real one.

We propose the validity ladder (Figure~\ref{fig:ladder}), which supplies a statistic, a pass rule and a human reference for each use, following the view in validity theory that each interpretation and use of a score needs its own evidence \citep{messick1995validity,kane2013validating,lin2026constructs}.
A gate measures how many repeated draws make a persona's score precise enough to read.
Above the gate, level 1 tests whether a single draw is an answer pattern a person could give, and level 2 tests whether gaps between groups match the population's.
Level 3 tests whether the population level matches, and level 4 tests whether the items relate to one another as they do in people.
A simulated sample supports only the claims whose rungs it has passed, and the rules are released as code.

In the worked example, the ladder is applied to PHQ-8 depression assessments that four language models produced as demographic personas \citep{keough2026plausible}, read against a national health survey.
Controls built from real survey respondents then test whether the ladder passes real people and catches failures planted at known sizes, a test any check on simulated samples has to pass before its verdicts on models carry weight.
The ladder is then applied to ten releases by other authors, ranging from silicon samples and digital twins to a model fine-tuned on survey answers \citep{suh2025subpop} and a personality inventory.
Together the three parts ask whether the ladder separates samples that stand in for people from samples that do not, on instruments and designs it was not built around.
The next section sets out each rung, the use it serves and what a study needs to run it.
